% BUBBLE-BLUE conference manuscript
% Overleaf: New Project → Blank → upload main.tex and refs.bib
% Compiler: pdfLaTeX  (Menu → Compiler)
% IEEEtran is provided by Overleaf.

\documentclass[conference,10pt]{IEEEtran}

\usepackage[T1]{fontenc}
\usepackage[utf8]{inputenc}
\usepackage{cite}
\usepackage{amsmath,amssymb}
\usepackage{booktabs}
\usepackage{tabularx}
\usepackage{array}
\usepackage{url}
\usepackage{hyperref}

\hyphenation{Blue-tooth BUBBLE-BLUE L2CAP}

\begin{document}

\title{Implementing BUBBLE-BLUE over BLE L2CAP:\\
An Open Mesh with Private Groups}

\author{
\IEEEauthorblockN{Paul-Maxence Baraton, Nadjib Achir, Philippe Jacquet}
\IEEEauthorblockA{INRIA, France\\
\{paul-maxence.baraton, nadjib.achir, philippe.jacquet\}@inria.fr}
}

\maketitle

\begin{abstract}
BUBBLE-BLUE (BB) is a multi-hop private group that rides on an \emph{open}
Bluetooth mesh: any BB node may forward, only holders of a group key may
read. Prior work specified the construction and evaluated Connected
Dominating Set (CDS) election in simulation. This paper reports the first
implementation, on BLE L2CAP Connection-Oriented Channels, without
Bluetooth pairing. We describe the on-wire format, the split between a
topological plane and an application plane, and a repeatable evaluation on
Qualcomm Linux boards (RUBIK Pi~3 and Arduino UNO~Q). A three-node mesh
completes in 7--14\,s, remains complete for 97--98\% of a steady run, and
delivers 100\% of chat. Driving the shipped forwarder in simulation, the
smallest TTL that reaches every node equals the source's hop eccentricity:
the CDS adds no path stretch. The practical limit is therefore a hop budget
(default TTL~10 spans about twenty nodes at mean degree 4--5). AES-256-GCM
runs at $30\text{k}$--$173\text{k}$ packets/s, far above the BLE link; the
Arduino host stack saturates near 10 messages/s. Revocation excludes the
victim in every trial and strands legitimate members only inside the join
window, before public keys have flooded.
\end{abstract}

\begin{IEEEkeywords}
Bluetooth Low Energy, L2CAP, mesh networks, connected dominating set,
group communication.
\end{IEEEkeywords}

\section{Introduction}

Infrastructure-less local communication still collapses, in practice, onto
a Wi-Fi access point or a smartphone OS that will not host a third-party
mesh. Bluetooth is ubiquitous and needs no operator. The SIG Mesh
profile~\cite{mesh-profile}, however, is a different stack, designed for
managed flooding (lighting, sensors), not for a private group that forms
between people who share a key and nothing else.

BUBBLE-BLUE~\cite{achir2025bb} proposed the opposite split: an
\emph{open topological mesh} that any BB node may join and relay, and
\emph{private application bubbles} readable only with a group key.
Routing is a dynamic CDS. That note gives the algorithms and unit-disk
simulations. It does not give a BLE transport, a frame layout, a
revocation protocol that reaches two-hop members, or measurements on
hardware.

This paper fills that gap. The implementation runs on two Qualcomm Linux
boards over BLE L2CAP CoC~\cite{btcore}. Connections use
\texttt{BT\_SECURITY\_LOW}: there is no SMP pairing, no bond, and no
address allow-list. Membership is the group key. A node that has never met
another can still talk to it. A passive sniffer in range can read routing
headers; it cannot read payloads (Section~\ref{sec:security}).

Our contributions are:
\begin{enumerate}
\item A complete BB stack on BLE CoC, including the 26-byte routing
  header, packet types, opaque forwarding, and epoch-based revocation
  (Section~\ref{sec:design}).
\item A repeatable multi-board evaluation method and the figures it
  produces on two platforms (Section~\ref{sec:eval}).
\item Characterised operating limits: hop-TTL, join-window revocation,
  and per-platform saturation---including two implementation defects the
  limit studies exposed and that are now fixed
  (Section~\ref{sec:defects}).
\end{enumerate}

The rest of this paper is organised as follows.
Section~\ref{sec:related} positions BB against CDS literature and existing
BLE messengers. Section~\ref{sec:design} presents the design.
Section~\ref{sec:eval} evaluates the implementation.
Section~\ref{sec:concl} concludes.

\section{Related work}
\label{sec:related}

\paragraph{CDS and multi-hop broadcast.}
Wu and Li's 1999 local rule~\cite{wuli1999} is our default election; their
2001 scheme~\cite{wuli2001} is the broader dominating-set routing setting.
Multipoint relays~\cite{qayyum2002} are the OLSR
construction~\cite{rfc3626}; in BB, MPR is an alternative
\emph{algorithm} whose \emph{result} is still a CDS. Jacquet's analytic
densities~\cite{jacquet2004} and the 1D/2D unit-disk comparison
in~\cite{achir2025bb} are why Wu--Li was preferred for a first
implementation: simpler, and better than MPR on the 2D disk in that note.
Dumbo~II~\cite{dumbo2008} is the emergency-network precedent cited
in~\cite{achir2025bb}. BLE has no native IP-style broadcast: every local
flood is a sequence of L2CAP unicasts. That cost model is why a CDS
(fewer transmitters) matters more here than on Wi-Fi MANETs.

\paragraph{BLE transports.}
We use LE Connection-Oriented Channels~\cite{btcore}, not RFCOMM
(Classic, serial, no BLE equivalent) and not GATT as the data plane
(serialised writes, no native fan-out to CDS neighbours). Advertising
carries \emph{presence} only. SIG Mesh~\cite{mesh-profile} remains a
separate, provisioned stack.

\paragraph{Adjacent messengers.}
BitChat~\cite{bitchat} shows that a BLE-only chat can exist at consumer
scale; it floods over GATT and elects no backbone. Briar~\cite{briar} is
delay-tolerant sync between \emph{pre-paired} contacts over Classic
RFCOMM, Wi-Fi and Tor---a different problem. Reticulum~\cite{reticulum} is
the closest Linux BLE stack; it is identity-routed, not CDS-routed.
Bridgefy and BLEMeshChat likewise flood or store-and-forward without a
CDS. None of these is a BB node. Netcore Causeway~\cite{causeway} shares
high-level aims with different technology; we do not treat it as a
baseline.

\section{Design}
\label{sec:design}

A BB node is a Linux process that owns an HCI adapter, accepts CoC
connections on PSM \texttt{0x0081}, and runs the control and forwarding
logic below. Pairing is unused. Confidentiality is application-layer.

\subsection{Open mesh, private bubble}

Two planes share the same links and the same 26-byte header, and nothing
else.

The \textbf{topological plane} is open. Any BB node may advertise
presence, accept a connection, take part in CDS election, and relay a
frame it cannot decrypt. AES-GCM failure is logged; it is never a reason
to drop the relay.

The \textbf{application plane} is private. Group traffic, key
announcement and revocation are readable only with the current group
key. Membership is that 32-byte secret, not a Bluetooth bond.

This is the construction of~\cite{achir2025bb}, realised as opaque
forwarding of group data by nodes that may not belong to the bubble.

\subsection{Presence and sparse links}

A node advertises a compact presence record in a connectable legacy
advertisement (manufacturer data, magic \texttt{BB}): anonymous
identifier, CDS flag, and the BD\_ADDR at which it accepts CoC. Scanners
accumulate a \emph{presence catalog} indexed by anon\_id. The catalog
survives a dropped socket and is how a neighbour is redialled without
pairing.

Links are not a full mesh. Each node caps its L2CAP degree and prefers
CDS successors. On a six-node soak, boards held two or three links each
and still reached the other five. When the CDS is empty---at start-up, and
whenever legacy advertisements cannot carry neighbour lists under a
$\sim$27-byte budget---a node dials from the catalog rather than waiting
for an air-only election that cannot see two-hop topology.

The implementation runs on RUBIK Pi~3 (Ubuntu, Bluetooth~5.0) and Arduino
UNO~Q (Debian plus an STM32, Bluetooth~5.1). CPython~3.14 is required for
LE 4-tuple sockets. A privileged daemon owns the adapter; the UI is an
unprivileged client, so restarting a display does not drop links.

\subsection{Routing header and packet types}

Every frame begins with a 26-byte cleartext header:
\begin{center}
\texttt{type(1) | src(6) | dst(6) | ttl(1) | seq(4 BE) | ts\_ms(8 BE)}
\end{center}
Relays decrement TTL on the wire, so TTL is \emph{not} in the AES-GCM
additional authenticated data; the rest of the header is. A relay without
the group key can still forward, and cannot forge an AAD-covered field.
Broadcast destination is \texttt{FF:FF:FF:FF:FF:FF}.

The ciphertext is AES-256-GCM~\cite{nist-gcm,rijndael} with a 16-byte
tag. There is no nonce on the wire: the per-packet key is derived from
the group or pairwise secret, source, sequence and timestamp. Overhead is
42 bytes. Against a conservative 512-byte MTU that leaves 470 bytes of
plaintext. Both boards negotiate L2CAP MTU 672, so a maximum-size
message is still one frame; 512 is policy. A 471-byte payload is refused,
never truncated.

\begin{table}[t]
\centering
\caption{BB packet types.}
\label{tab:pkts}
\begin{tabular}{@{}llll@{}}
\toprule
Type & Name & Plane & Role \\
\midrule
\texttt{0x01} & HELLO & both & Neighbour table, epoch, CDS \\
\texttt{0x02} & DATA & private & Group broadcast \\
\texttt{0x04} & SUBGROUP & private & Addressed multicast \\
\texttt{0x05} & REKEY & private & Pairwise epoch catch-up \\
\texttt{0x06} & ANNOUNCE & private & Signed pubkey binding \\
\texttt{0x07} & REKEY\_FLOOD & private & Coordinator revocation \\
\bottomrule
\end{tabular}
\end{table}

Two sequence counters number these frames (Hello/ANNOUNCE versus
DATA/REKEY). They must not share an anti-replay window
(Section~\ref{sec:defects}).

\subsection{CDS, forwarding, and membership expiry}

Roles follow Wu--Li 1999~\cite{wuli1999}. MPR~\cite{qayyum2002} is
implemented in the simulator; boards run Wu--Li. A node forwards if it is
RELAY, or if it has two or more L2CAP links (a bridge the CDS has not
named---the usual case at the edge of range, where the advertisement
graph is a clique and the L2CAP graph is a star). Duplicate suppression
is keyed by $(\mathit{src},\mathit{seq})$ and by which counter produced
$\mathit{seq}$.

A neighbour is dropped after \texttt{HELLO\_MISSES} consecutive silent
Hello periods (default 3, minimum 2), i.e.\ 15\,s at a 5\,s interval.
Counting periods keeps the threshold meaningful if the interval is
retuned.

\subsection{Security and revocation}
\label{sec:security}

Link encryption is off. Group AEAD is AES-256-GCM under the group key;
pairwise keys are X25519~\cite{rfc7748} after ANNOUNCE; rekeys are
attested with Ed25519~\cite{rfc8032}. The routing header is visible to a
sniffer; payloads are not.

The coordinator is the member with the smallest anon\_id among
\emph{all known members} (Hello-resolved \emph{and} ANNOUNCE-learned).
Election and validation of a rekey must use that same set
(Section~\ref{sec:defects}).

Only the coordinator may revoke. It bumps the epoch, rotates the group
key, and floods \texttt{REKEY\_FLOOD} wrapped under \emph{pre-bump}
pairwise keys, excluding the victim. Late members catch up on the next
stale Hello. The coordinator cannot wrap a key it does not hold: Hello
supplies direct neighbours only; everyone else arrives in ANNOUNCE
(30\,s). Revoking inside that \emph{join window} is refused or strands
legitimate members---an operating limit, not a demonstration surprise.

Both SoCs use ARMv8 Cryptography Extensions (\texttt{gcm-aes-ce} via
OpenSSL). Qualcomm QCE is present but does not register GCM; it is not on
the path.

\section{Evaluation}
\label{sec:eval}

Experiments use simultaneous probes on all boards, the same git commit,
and \emph{worst-board} figures (a mesh is complete when the last node has
joined). Warm-up (25\,s) is counted as convergence, not as instability.
Departure is timed on the presence table, not on L2CAP sockets (a socket
dies with its process regardless of \texttt{HELLO\_MISSES}). Load tests
use a single sender; each message carries an index, and delivery is that
index seen by a listener. Simulation instantiates the shipped forwarder
on a unit-disk radio and drains floods breadth-first (call-stack,
depth-first delivery is not a radio).

\subsection{Convergence, stability, and departure}

\begin{table}[t]
\centering
\caption{Mesh formation and chat delivery.}
\label{tab:conv}
\begin{tabular}{@{}lcccc@{}}
\toprule
Setting & $n$ & Full mesh & Steady complete & Delivery \\
\midrule
Automated & 3 & 7--14\,s & 97--98\% & 100\% \\
Automated & 2 & 3--14\,s & 97\% & 100\% \\
Manual soak & 6 & $\sim$40\,s & stable after 60\,s & 427/429 \\
\bottomrule
\end{tabular}
\end{table}

The two losses in 429 were the first message (before convergence) and the
last (during shutdown). Applications should display mesh state rather
than assume delivery in the first tens of seconds. Relay role flips in
the automated runs: 0. Steady-state L2CAP flaps: a few per run, all
recovered.

\begin{table}[t]
\centering
\caption{Departure delay vs.\ configured Hello misses (presence view).}
\label{tab:misses}
\begin{tabular}{@{}cccc@{}}
\toprule
\texttt{HELLO\_MISSES} & Expected & Measured departure & Rejoin \\
\midrule
3 & 15\,s & 17--21\,s & 7--12\,s \\
6 & 30\,s & 32--36\,s & 7--10\,s \\
\bottomrule
\end{tabular}
\end{table}

Killing a node without notice (Table~\ref{tab:misses}): departure tracks
the configured miss count. Rejoin does not: it is discovery and dial.

\subsection{Crypto is not the bottleneck}

\begin{table}[t]
\centering
\caption{AES-256-GCM on a 470-byte payload (mean of 20\,000 calls).}
\label{tab:crypto}
\begin{tabular}{@{}lcc@{}}
\toprule
Platform & Encrypt + decrypt & Equivalent rate \\
\midrule
RUBIK Pi~3 & $3.0+2.8\,\mu$s & $173\,000$ pkt/s \\
Arduino UNO~Q & $18.2+15.2\,\mu$s & $30\,000$ pkt/s \\
\bottomrule
\end{tabular}
\end{table}

BLE CoC carries a few hundred packets per second. AES-GCM is two to three
orders of magnitude above the link (Table~\ref{tab:crypto}). Energy of
the crypto path is not measured.

\subsection{Topological and revocation limits}

At mean degree $\sim$4--5, the minimum TTL that delivers to every node
equals the source eccentricity on every draw (chains and random geometric
graphs):

\begin{table}[t]
\centering
\caption{Minimum TTL for full delivery equals hop diameter
  (Wu--Li, shipped forwarder, geometric radio).}
\label{tab:ttl}
\begin{tabular}{@{}rrrrr@{}}
\toprule
$n$ & Degree & Diam.\ (hops) & Min TTL & Deliv.\ @ TTL 10 \\
\midrule
10 & 3.9 & 4 & 4 & 100\% \\
20 & 4.3 & 8 & 8 & 100\% \\
30 & 3.3 & 11 & 11 & 93\% \\
50 & 5.2 & 13 & 13 & 94\% \\
100 & 5.2 & 19 & 19 & 23\% \\
\bottomrule
\end{tabular}
\end{table}

There is no ``it breaks at $N$ nodes''. Past a 10-hop diameter, default
TTL~10 no longer spans the network---about twenty nodes at this density.
At six nodes the default is ample. Below mean degree $\sim$5, no connected
geometric draw in 100 trials (percolation). Denser graphs elect fewer
relays (52\%$\to$38\%) and cost more per flood.

After ANNOUNCE has run, revocation at $n=6$--$30$: 0 stranded members, 0
leaks to the victim. Cost grows from 12 to 83 emissions per node
(roughly quadratic). Before ANNOUNCE, revocation is often refused.
Rule: wait one ANNOUNCE period after a join.

\subsection{Load and two defects the study caught}
\label{sec:defects}

A single sender, 64--470 byte payloads. The RUBIK listener received
100\% through 30\,msg/s. The UNO~Q tracked to 10\,msg/s, lost at 15, then
delivered 0---including a 1\,msg/s recovery step after 25\,s idle---while
its links stayed nominally up. The kernel logged
\texttt{Unable to allocate ident: -28} (\texttt{ENOSPC} on BlueZ L2CAP
\emph{signalling} identifiers); controller error counters stayed 0.
Data-plane saturation delays Hellos, neighbours expire, auto-dial
bursts, the ident pool empties. Restarting the Bluetooth service recovers
the board. A separate QCA firmware stall after $\sim$6 weeks of uptime
requires a reboot; it is not this failure.

Limit studies also found two protocol bugs, both now fixed.

\paragraph{Shared anti-replay window.}
Hello/ANNOUNCE and DATA were numbered from independent counters but
checked against one window. A DATA frame whose sequence equalled an
earlier ANNOUNCE from the same node was dropped by every receiver and
every relay. Reproduction: sequence 4 delivered 5/5, sequence 5 (the
ANNOUNCE) 0/5, sequence 6 delivered 5/5. Likely cause of the 2/429 soak
losses. Fix: one window per counter.

\paragraph{Split-brain coordinator.}
Election used Hello-resolved neighbours only; validation of a rekey also
counted ANNOUNCE-learned peers. On a stringy 30-node graph (mean
degree~3.3) four nodes declared themselves coordinator; revocation was
refused mesh-wide while the initiator's epoch advanced; 28 of 29
legitimate members stayed on the old key. Fix: one membership view for
both.

\section{Conclusion}
\label{sec:concl}

An open BLE mesh with private groups is implementable without pairing.
On the boards we have, the mesh converges in seconds, the CDS does not
stretch paths, and the default TTL is a hop budget rather than a magic
node count. Revocation works once keys have been announced and fails in
a documented window if they have not. Cryptography is not the throughput
limit; the Arduino host stack is.

Physical range, indoor multi-hop lines, and walking partition/merge still
require the pending displays (an operator must read peer count without
SSH). Those experiments are specified; they are not in this paper.
Energy of AES-GCM versus the radio, a six-node repeat of
Section~\ref{sec:eval}, and an Android CoC port are future work.

\section*{Acknowledgment}

The authors thank Wilfred (Qualcomm) for platform guidance and for the
questions that produced the accelerator and host-stack measurements. The
protocol note~\cite{achir2025bb} acknowledges Matthieu Finiasz (OLVID)
and Thomas Jaussouin and Olivier Berni (Lunabee); that debt carries over.

\bibliographystyle{IEEEtran}
\bibliography{refs}

\end{document}
